\section{Export and Runtime}
\label{sec:deployment}
\subsection{Exported Graph and Strict Compilation}
The retained ONNX artifact uses opset 20 and a static $1\times3\times640\times640$
input. Inspection finds 89 convolutions, 76 sigmoid and 76 multiplication nodes,
9 splits, 14 additions, 20 concatenations, 3 max-pooling operations, 2 resizes,
and one softmax. There are no MatMul, Gemm, Reshape, Transpose, Div, or Slice
operators. The sigmoid/multiplication pairs represent retained SiLU activations;
operator counts before fusion are not execution-kernel counts.

Compilation targets DLA core 0 in FP16 with GPU fallback disabled and explicit
native I/O formats. Successful compilation is checked against both the verbose
placement report and the optimized-engine inspector. A one-node DLA inspector
result is evidence of a single compiled loadable, not a claim that the original
network contains one operation. The build and profile records remain linked to
the engine name, compiler version, and output bindings.

This audit is configuration-specific. TensorRT's DLA softmax support depends on
the target and shape, and the compiler warns that the selected approximation
can introduce numerical error \cite{nvidia_dla_working}. A source-level
rewrite that compiles on this Orin/TensorRT combination is not sufficient
evidence for another DLA generation, input size, or precision. Build success
also does not validate raw-output decoding.

\subsection{Native Layouts and the Host Boundary}
\label{sec:packing}
The strict build requests \texttt{fp16:dla\_hwc4} input and
\texttt{fp16:chw16} outputs. Its inspector describes a DLA image input format
and sixteen-channel-vectorized outputs. For this three-channel image, the
physical input buffer pads channels to four and uses interleaved pixel storage;
the $640$-pixel row satisfies the documented 64-byte alignment. The buffer
occupies $4\times640\times640\times2=3{,}276{,}800$ bytes. Each output has
144 channels, already divisible by 16, so CHW16 requires no extra channel
padding.

The reference runner converts between these buffers and logical NCHW tensors
on the CPU using tensor views, permutations, materialization, and copies.
Its timed stages include both packing and H2D transfer at the input, and both
D2H transfer and unpacking at the output. They cannot be interpreted as pure
layout-conversion timings. A CUDA stream and device-visible buffers are still
used for execution and transfer even when the learned graph has no GPU layers.

Logical shape is insufficient to allocate or interpret a native binding.
Runtime validation must record format descriptions, vectorized dimensions,
components per element, and context strides, then test the corresponding
packing rules. The retained summary uses a CHW4 shorthand for the input;
we follow the build's requested DLA format and require a runtime record before
treating format names as interchangeable.
\missing{P4}{archive binding formats/strides and raw multi-image timing records}.

\subsection{Confidence-First Host Decoding}
\label{sec:decode}
The host splits each $P_\ell$ into 64 regression channels and 80 class channels.
For confidence threshold $\tau\in(0,1)$, an anchor survives if
\begin{equation}
 \max_c S_{\ell,c}(h,w) > \log\frac{\tau}{1-\tau}.
 \label{eq:threshold}
\end{equation}
This is equivalent in exact arithmetic to applying sigmoid and then the same
maximum-class threshold. Regression decoding is independent of class
confidence, so it can be postponed until after this rejection step. For each
survivor and box side $q\in\{l,t,r,b\}$, the DFL expectation is
\begin{equation}
 \hat d_q=\sum_{j=0}^{r-1}j\,
 \frac{\exp R_{qj}}{\sum_{k=0}^{r-1}\exp R_{qk}}.
 \label{eq:dfl}
\end{equation}
The decoder places cell centers at $(w+0.5,h+0.5)$, forms the four box
coordinates from these distances, scales by the level stride, applies class
sigmoid, and performs class-aware NMS. The final coordinates are mapped back
through the image's letterbox transform.

If $m$ of $N=8400$ anchors survive, dense class scanning remains
$O(Nn_c)$, while the regression softmax/expectation work falls from $O(4rN)$
to $O(4rm)$. This does not reduce raw-output transfer volume. The repeated-image
summary reports $m=13$ at $\tau=0.25$, a favorable case that cannot establish
the candidate distribution on a dataset.

The equivalence argument assumes identical threshold strictness, class
selection, NMS settings, and tie handling. Floating-point behavior near the
threshold, multi-label validation, and NMS ordering require explicit tests.
\missing{P3}{measure raw-output agreement and dense/sparse decoder parity,
including low-confidence dataset validation}.
